% ECONOMICS_PROBLEM: {"id": "D72-197", "title": "Democracy's economic mechanisms and resilience", "jel_code": "D72", "source_id": "OP-0388"}
% !TeX program = xelatex
\documentclass[11pt,a4paper]{article}
\usepackage[margin=23mm]{geometry}
\usepackage{fontspec}
\setmainfont{Latin Modern Roman}
\usepackage{amsmath,amssymb,mathtools}
\usepackage{xcolor,enumitem,booktabs,longtable,array,xurl,hyperref}
\definecolor{muted}{HTML}{53636C}
\hypersetup{colorlinks=true,urlcolor=blue,linkcolor=blue}
\setlength{\parindent}{0pt}
\setlength{\parskip}{5pt}
\newcommand{\R}{\mathbb R}
\newcommand{\E}{\mathbb E}
\newcommand{\N}{\mathbb N}
\newcommand{\ind}{\mathbf 1}
\newcommand{\Var}{\operatorname{Var}}
\newcommand{\Cov}{\operatorname{Cov}}
\newcommand{\supp}{\operatorname{supp}}
\DeclareMathOperator*{\argmin}{arg\,min}
\DeclareMathOperator*{\argmax}{arg\,max}
\newcommand{\entrymeta}[3]{{\sffamily\small\color{muted}\textbf{\texttt{#1}}\quad|\quad #2\par #3\par}}
\newcommand{\Problemref}[1]{\texttt{#1}}
\newcommand{\JELref}[2]{\texttt{#2} (JEL \texttt{#1})}
\begin{document}
\section*{Democracy's economic mechanisms and resilience}
\textbf{Problem ID:} \texttt{D72-197}\quad \textbf{JEL:} \texttt{D72}\par
% BEGIN ECONOMICS STATEMENT
\label{problem:OP-0388}
{\sffamily\small\color{muted}Original subject group: Political economy and institutions\par}
\label{definitions:problem:30007670}
\textbf{Audit classification:} Background; proposed extension. Original 2021 NBER metadata and revised primary manuscript checked. It already estimates exposure-support links and suggestive resilience; the broader feedback-loop specification is editorial.
\subsection*{Definitions and precise research target}
\textbf{Two distinct causal targets (editorial).} Fix baseline countries, a finite horizon \(T\), a democratic-status coding rule, normalized country weights, and fixed citizen cohorts. Democratic status \(D_{ct}\in\{0,1\}\) is based on prespecified electoral contestability, participation, and executive constraints. Real income per resident \(Y_{cT}\) uses a common price base and a declared resident denominator. Procedural support \(S_{icT}\) is a measured survey index distinct from approval of the current incumbent. All relevant outcomes have finite means, and citizen weights \(\mu_{ic}\) are nonnegative and sum to one within country.

A controlled exposure path \(d=(d_0,\ldots,d_T)\) assigns democratic status at every date under a stated intervention definition. It can define income and support contrasts
\[
 \tau_Y^d=E[Y_{cT}(d^1)-Y_{cT}(d^0)],\qquad
 \tau_S^d=E\sum_i\mu_{ic}[S_{icT}(d^1)-S_{icT}(d^0)].
\]
Such paths must specify timing, anticipation, financing, and the political rules being changed. A mere change of a regime label is not an intervention. Full status-path assignment fixes democratic survival by construction, so it is unsuitable for a nontrivial resilience target.

For resilience, instead fix countries democratic at baseline \(t_0<T\), and compare feasible economic, information, or institutional reform regimes \(p\) and \(p_0\) that preserve the initial-status requirement while allowing future status to respond endogenously. Each regime specifies implementation, timing, public information, and resource use; it does not set future GDP, citizen support, or democratic status directly. Let \(D_{ct}(p)\), \(Y_{cT}(p)\), and \(S_{icT}(p)\) be the resulting outcomes, and define
\[
 Z_{cT}(p)=\mathbf 1\{D_{ct}(p)=1\text{ for every }t=t_0,\ldots,T\},\qquad
 \tau_Z^p=E[Z_{cT}(p)-Z_{cT}(p_0)].
\]
This survival indicator is an outcome, not an assigned path. Define analogous regime-specific income and procedural-support contrasts. Never combine a controlled-path contrast and an endogenous-survival contrast as if they arose from the same intervention.

The task is to characterize the identified sets for these contrasts, and determine which economic or institutional mechanisms connect democratization, prosperity, procedural support, and durability. Declare the admissible causal model class, transition/reform selection, global-shock and interference rules, attrition and cohort measurement, and any mediation assumptions. Transition timing is not randomized. The average income effect of democratization, a correlation between income and support, and a causal effect of truthful disclosure are different quantities. Randomized economic or information interventions can identify parts of the feedback chain but do not identify unrestricted mediation without additional exclusions and cross-world or sequential assumptions. The cited paper already supplies democratic-exposure/support estimates and suggestive resilience evidence; the combined problem here is a proposed extension, with no verified claim that this exact formulation is a published open conjecture.
\subsection*{Variants and assumption changes}
\begin{enumerate}
\item Exposure and beliefs (source-informed editorial): use cohorts with distinct democratic exposure and independently varied performance information, separating actual experience from perceptions under explicit exclusion restrictions.
\item Shock resilience (editorial extension): define a prospective set of economic shocks and compare institutional-survival responses under alternative accountability institutions. Test whether baseline procedural support predicts survival and separately identify any causal support intervention.
\end{enumerate}
\subsection*{References and source audit}
\textbf{Support and access.} Original 2021 NBER metadata and revised primary manuscript checked. It already estimates exposure-support links and suggestive resilience; the broader feedback-loop specification is editorial. \textbf{Locator:} MIT author manuscript, October 2, 2023, abstract. \textbf{Versions:}  
\begin{enumerate}\raggedright
\item\hypertarget{definitions:source:30007670:1}{} Daron Acemoglu; Nicolás Ajzenman; Cevat Giray Aksoy; Martin Fiszbein; Carlos A. Molina. 2021. \emph{(Successful) Democracies Breed Their Own Support}. Abstract, paragraphs on democracy exposure, economic performance, and support.
\url{https://www.nber.org/papers/w29167}.
DOI: \url{https://doi.org/10.3386/w29167}.
\item\hypertarget{definitions:source:30007670:2}{} Daron Acemoglu; Nicolás Ajzenman; Cevat Giray Aksoy; Martin Fiszbein; Carlos Molina. 2023. \emph{(Successful) Democracies Breed Their Own Support}. October 2, 2023 author manuscript, abstract.
\url{https://economics.mit.edu/sites/default/files/inline-files/SuccessfulDemocracies2023.pdf}.
\end{enumerate}

\subsection*{Common conventions: Source mathematical conventions}

These conventions apply to each entry unless explicitly replaced.

\textbf{Models and identification.} A primitive model \(m\) specifies all probability laws, feasible choices, information, equilibrium and measurement rules used by its target. The admissible class \(\mathcal M\) is fixed independently of outcomes. If \(P_O(m)\) is its observable probability law and \(\tau(m)\) the target, its identified set at observable law \(P\) is
\[
 \mathcal I_\tau(P)=\{\tau(m):m\in\mathcal M,\ P_O(m)=P\}.
\]
Point identification means that this set is a singleton. An empty set signals incompatibility of data and assumptions. Coordinate bounds need not describe the joint set. Bounds are sharp only if every retained target is attainable or arbitrarily approachable by compatible models. Finite bounds require explicit restrictions on unobserved quantities; unrestricted confounding, hidden wealth, extrapolation or arbitrary equilibrium selection can yield uninformative sets.

\textbf{Causal interventions.} A potential outcome \(Y_i(p)\) is the outcome under a complete policy regime \(p\), including timing, financing, information and market spillovers. The target population and units are fixed before treatment. With interference, use the complete assignment vector or a justified exposure mapping; individual no-interference notation alone cannot identify an economy-wide intervention. A difference of mean potential outcomes does not require the cross-world joint distribution. Conditioning on post-treatment migration, survival, enrollment or employment changes the estimand and needs separate justification.

\textbf{Statistical statements.} An estimator, confidence set, forecast or test specifies a sampling law, model class and loss/coverage criterion. Uniform \((1-\alpha)\)-coverage means \(\inf_{m\in\mathcal M}\Pr_m\{\tau(m)\in C_n\}\geq1-\alpha\), or an explicitly asymptotic version. The whole parameter space trivially has coverage, so informativeness requires a stated width, risk or power objective. A forecasting improvement is relative to a named benchmark, information set and evaluation population.

\textbf{Equilibrium and optimization.} An equilibrium comprises feasible plans, optimality, beliefs where needed, budget constraints and all market/resource clearing. The primitives or a stated benchmark define each correspondence \(\mathcal E(p;m)\). Nonemptiness is assumed or proved, never inferred just from compact policy sets. A selected-equilibrium target uses a declared selection rule; a robust target quantifies over all permitted equilibria. Use suprema or infima unless attainment is established. An \(\varepsilon\)-optimal policy is feasible and within \(\varepsilon\) of the finite optimum value. Report an empty feasible set separately.

\textbf{Welfare and accounting.} Normative utility, interpersonal normalization, social weights, numeraire, discounting and the use of public receipts are primitives. Behavioral utility need not coincide with the welfare criterion. Household welfare includes ownership income and rents once; adding profits again to compensating variation generally double-counts them. A monetary equivalent is defined by an explicit transfer and price convention, with monotonicity and range conditions. Average effects, mechanism contributions, transfers and welfare losses are different targets. Infinite sums require convergence and dynamic feasible plans require terminal or no-Ponzi conditions.

\textbf{Source versions.} A working-paper date, revision date and journal publication year are distinguished. A DOI for a journal version is not automatically the DOI of an earlier manuscript. Locators refer to the stated version and printed pages where specified. A metadata-only check does not verify a claim about a full-text conclusion. Every entry's source subsection narrows the provenance claim accordingly.
% END ECONOMICS STATEMENT
\end{document}
